\section{Encounter-quotient visibility}

Let
\[
E=\{1,a,b,ab\}\cong C_2\times C_2
\]
and define the natural encounter quotient
\[
q:\langle A,B\rangle\longrightarrow E,
\qquad
q(A)=a,
\qquad
q(B)=b.
\]
It follows that
\[
q(H)=ab,
\qquad
q(\comm)=q(H^2)=1.
\]
The kernel is exactly
\[
\Ker(q)=\langle H^2\rangle=\langle\comm\rangle,
\]
which has order (n).

The binary return receipt has image
\[
q(\receipt)=q(H^n)=(ab)^n.
\]
Therefore
\[
q(\receipt)=
\begin{cases}
1, & n\text{ even},\\
ab, & n\text{ odd}.
\end{cases}
\]

\begin{proposition}[Parity-sensitive concealment]
The encounter quotient hides the binary return receipt if and only if the
visible cycle length (n) is even.
\end{proposition}

The same parity law has now appeared in two equivalent forms:
\[
\receipt\in\langle\comm\rangle
\quad\Longleftrightarrow\quad
q(\receipt)=1
\quad\Longleftrightarrow\quad
n\text{ is even}.
\]

